\hypertarget{c11-standard-library-additions}{%
\chapter{C++11 STANDARD LIBRARY
ADDITIONS}\label{c11-standard-library-additions}}

C++11 massively expanded the STL.

\begin{center}\rule{0.5\linewidth}{0.5pt}\end{center}

\hypertarget{unordered-containers-hash-maps}{%
\section{1. UNORDERED CONTAINERS (HASH
MAPS)}\label{unordered-containers-hash-maps}}

C++98 \passthrough{\lstinline!std::map!} is a tree (O(log n)). C++11
added hash tables (O(1) average).

\begin{itemize}
\tightlist
\item
  \passthrough{\lstinline!std::unordered\_map!}
\item
  \passthrough{\lstinline!std::unordered\_set!}
\item
  \passthrough{\lstinline!std::unordered\_multimap!}
\item
  \passthrough{\lstinline!std::unordered\_multiset!}
\end{itemize}

\begin{lstlisting}[language={C++}]
#include <unordered_map>
std::unordered_map<std::string, int> scores;
scores["Alice"] = 100; // O(1)
\end{lstlisting}

They require a Hash function for the key type.

\begin{center}\rule{0.5\linewidth}{0.5pt}\end{center}

\hypertarget{stdarray}{%
\section{2. STD::ARRAY}\label{stdarray}}

\passthrough{\lstinline!std::array!} is a fixed-size wrapper around
C-style arrays. It doesn't decay to a pointer automatically and knows
its size.

\begin{lstlisting}[language={C++}]
#include <array>
std::array<int, 5> arr = {1, 2, 3, 4, 5};
// arr.size() is 5
\end{lstlisting}

Prefer over C-arrays (\passthrough{\lstinline!int arr[5]!}).

\begin{center}\rule{0.5\linewidth}{0.5pt}\end{center}

\hypertarget{stdtuple}{%
\section{3. STD::TUPLE}\label{stdtuple}}

Generalization of \passthrough{\lstinline!std::pair!} for N elements.

\begin{lstlisting}[language={C++}]
#include <tuple>
auto t = std::make_tuple(10, "Hello", 3.14);
int i = std::get<0>(t);
\end{lstlisting}

\begin{center}\rule{0.5\linewidth}{0.5pt}\end{center}

\hypertarget{regular-expressions}{%
\section{4. REGULAR EXPRESSIONS}\label{regular-expressions}}

\passthrough{\lstinline!std::regex!} provides regex matching and
replacement.

\begin{lstlisting}[language={C++}]
#include <regex>
std::regex pattern(R"(\d+)"); // Matches digits
bool match = std::regex_match("123", pattern);
\end{lstlisting}

\begin{center}\rule{0.5\linewidth}{0.5pt}\end{center}

\hypertarget{chrono-time}{%
\section{5. CHRONO (TIME)}\label{chrono-time}}

Type-safe time library.

\begin{lstlisting}[language={C++}]
#include <chrono>
auto start = std::chrono::high_resolution_clock::now();
// ... work ...
auto end = std::chrono::high_resolution_clock::now();
auto duration = std::chrono::duration_cast<std::chrono::milliseconds>(end - start);
\end{lstlisting}

\begin{center}\rule{0.5\linewidth}{0.5pt}\end{center}

\hypertarget{random}{%
\section{6. RANDOM}\label{random}}

Better random number generation than \passthrough{\lstinline!rand()!}.

\begin{lstlisting}[language={C++}]
#include <random>
std::random_device rd;
std::mt19937 gen(rd());
std::uniform_int_distribution<> dis(1, 6); // Dice roll
int roll = dis(gen);
\end{lstlisting}
